\scanpage{30}
\textbf{C. R. Acad. Sc. Paris, t. 292 (19 janvier 1981)}\hfill\textbf{Série I -- \mspage{215}}
\figure{\begin{tikzpicture}\draw (0,0)--(14,0);\end{tikzpicture}}

\textsc{Théorème} 2 (cf. [3]). -- \textit{Soit $m\in\mathbf N^*$ et soit $\mathrm{Mon}(\mathbf R^m)$ le monoïde des applications continues et injectives de $\mathbf R^m$ en $\mathbf R^m$. Si $\Omega$ est un faisceau sur $\mathcal V_m$ [au sens de l'axiome ($a$)] les conditions ci-dessous sont équivalentes :}

(i) \textit{$\Omega$ est localement constant} (cf. [4]);

(ii) \textit{$\Omega$ est séparé [au sens de l'axiome ($b$)];}

(iii) \textit{Le monoïde $\mathrm{Mon}(\mathbf R^m)$ opère sur $\Omega(\mathbf R^m)$ par automorphismes.}

Soit $\mathrm{L.\ C.}(\mathcal V_m)$ la catégorie des faisceaux localement constants sur $\mathcal V_m$ et soit $\Pi(m)$ le groupe symétrisé du monoïde $\mathrm{Mon}(\mathbf R^m)$.

\textsc{Théorème} 3 (cf. [3]). -- \textit{Le foncteur} [cf. \textit{théorème} 2, (iii)] :
\[
\Omega\to\Omega(\mathbf R^m):\quad\diagram{\begin{tikzcd}\mathrm{L.\ C.}(\mathcal V_m)\arrow[r]&\Pi(m)-\mathrm{Ens},\end{tikzcd}}
\]
\textit{est une équivalence de catégories} [\textit{i. e.} $\Pi_1(\widetilde{\mathcal V}_m)\xrightarrow{\sim}\Pi(m)$].

La théorie du degré de Brouwer fourni un homomorphisme de groupes :
\[
\deg_m:\quad\diagram{\begin{tikzcd}\Pi(m)\arrow[r]&\mu_2.\end{tikzcd}}
\]

\textsc{Théorème} 4 (cf. [3]). -- \textit{Les conditions ci-après sont équivalentes :}

($a$) \textit{l'homomorphisme $\deg_m:\Pi(m)\xrightarrow{\sim}\mu_2$ est un isomorphisme;}

($b$) \textit{tout homéomorphisme de $\mathbf R^m$ de degré $+1$ est « stable » (i. e. composé fini d'homéomorphismes dont chacun est l'identité sur un ouvert non vide).}

Puisque Kirby a prouvé que pour $m\ne4$ tout homéomorphisme de $\mathbf R^m$ de degré $+1$ est stable (\textit{cf.} [2]) on a le résultat suivant :

\textsc{Corollaire}. -- \textit{Si $m\ne4$ l'homomorphisme $\deg_m:\ :\Pi(m)\xrightarrow{\sim}\mu_2$ est un isomorphisme.}

\textsc{Définition}. -- On appelle \textit{triplet de Stokes} (en dimension $m$) tout triplet $\tau=(\Omega_m,\Omega_{m-1},\delta)$ d'un faisceau séparé $\Omega_m$ sur $\mathcal V_m$, d'un faisceau séparé $\Omega_{m-1}$ sur $\mathcal V_{m-1}$ et d'un morphisme fonctoriel :
\[
\delta:\quad\diagram{\begin{tikzcd}\Omega_m^\bullet\arrow[r]&\Omega_{m-1}\circ\partial\end{tikzcd}}
\]
($\partial:\mathcal V_m\to\mathcal V_{m-1}$ est le foncteur bord) qui soit un « isomorphisme local » au sens de l'axiome ($c$).

La notion de morphisme entre triplets de Stokes en dimension $m$ étant évidente, on note par $\mathrm{St}(m)$ la catégorie des triplets de Stokes en dimension $m$.

\textsc{Théorème} 5 (cf. [3]. -- \textit{Le foncteur} :
\[
(\Omega_m,\Omega_{m-1},\delta)\to\Omega_m(\mathbf R^m):\ \diagram{\begin{tikzcd}\mathrm{St}(m)\arrow[r,"\approx"]&\Pi(m)-\mathrm{Ens},\end{tikzcd}}
\]
\textit{est une équivalence de catégories.}

(*) Remise le 5 janvier 1981.

[1] R. \textsc{Godement}, \textit{Théorie des faisceaux}, Herman, Paris, 1958.

[2] R. \textsc{Kirby}, \textit{Ann. of Math.}, 89, 1969.

[3] M. \textsc{Wanderley}, \textit{Théorie axiomatique de l'orientation des variétés topologiques à bord (Thèse de} $3^{\mathrm e}$ \textit{cycle}, Montpellier, 1980).

[4] SGA 4, $1^{\mathrm{er}}$ vol., \textit{Séminaire de Géométrie algébrique de l'I.H.E.S.}, dirigé par M. \textsc{Artin}, A. \textsc{Grothendieck} et J. L. \textsc{Verdier} (\textit{Lecture Notes in Math.} 289, Springer).

\textit{U.E.R. de Mathématiques, Université des Sciences et Techniques du Languedoc;}\\
\textit{Ed. N. S. das Gracas. App. 202. Rua das Gracas. 277. 50.000 Recife. P. E., Brésil.}

C. R., 1981, $1^{\mathrm{er}}$ \textit{Semestre}. (T. 292, N° 3)\hfill\textbf{Série I -- 15}
\margin{29}
